\noindent
Transcranial focused ultrasound stimulation (tFUS) has emerged as a noninvasive neuromodulation approach with millimeter-scale focality and the ability to access subcortical and medial structures that are challenging to reach with electromagnetic methods such as transcranial magnetic stimulation (TMS) or transcranial direct current stimulation (tDCS) \cite{Naor2016-wz,blackmore2019mechanisms}. Foundational studies spanning \textit{in vitro}, small-animal, and human preparations established feasibility and provided early demonstrations of state-dependent effects \cite{Tyler2008-zs,Tufail2010-jd,legon2014transcranial}. These advances have catalyzed a growing body of research applying tFUS to probe and perturb human brain networks \cite{Legon2014-li,lee2016transcranial,Legon2018-ad,Legon2018-kb,sanguinetti2020transcranial,zhang2021transcranial,yaakub2023transcranial,lord2024transcranial}.

Despite this progress, the biophysical mechanisms by which low-intensity ultrasound influences neural activity remain an active area of investigation \cite{Blackmore2019-zo,kamimura2020ultrasound,dell2022current}. Proposed accounts include direct effects on mechanosensitive ion channels \cite{Sorum2021,Kubanek2018-eo,Yoo2022}, modulation of membrane capacitance and bilayer elasticity consistent with intramembrane cavitation \cite{krasovitski2011intramembrane,Plaksin2014}, flexoelectric coupling of lipid membranes \cite{Petrov2002}, and indirect astrocytic or neurovascular contributions \cite{mishima2024repetitive}. Clarifying these mechanisms is essential for developing principled strategies to bias excitation–inhibition balance and to achieve reproducible, directionally specific changes in circuit activity.

Concurrently, there is growing interest in how tFUS alters large-scale functional organization of the human brain. Functional connectivity (FC) -- typically operationalized as correlations in the blood-oxygen-level–dependent (BOLD) signal between regions -- provides a macroscopic readout that is sensitive to behavioral state and clinical phenotype \cite{rogers2007assessing,van2010exploring}. Human functional magnetic resonance imaging (fMRI) studies have reported both increases and decreases in FC following tFUS, likely reflecting differences in targeted circuits, sonication parameters, behavioral context, and analysis choices \cite{chou2024transcranial,yaakub2023transcranial,todd2018focused,meng2019resting,sanguinetti2020transcranial,zhang2021low,kuhn2023transcranial,lord2024transcranial}. A consistent picture has yet to emerge regarding when and where tFUS strengthens versus weakens coupling.

Here we focus on the subgenual anterior cingulate cortex (sgACC), a deep medial prefrontal node implicated in affect regulation and prominently linked to mood disorders \cite{drevets2008subgenual,greicius2007resting}. We conducted a within-subject, randomized, sham-controlled study in healthy adults (two sessions per participant) combining tFUS with concurrent fMRI. Our primary endpoint was BOLD FC between the sgACC and the rest of the brain, estimated at the subject level and analyzed with linear mixed-effects models.

We found that: (i) sgACC–whole-brain FC increased from pre- to post-sonication in Active relative to Sham (time$\times$condition Pre$\rightarrow$Post contrast), with a similar but non-significant trend during tFUS, (ii) this effect was not explained by regression-to-the-mean or baseline differences, as it persisted in baseline-controlled models, and (iii) most importantly, the added coupling was expressed selectively across canonical brain networks: Sham sessions showed an increase in sgACC coupling with the default mode network (DMN) over time, whereas Active tFUS preferentially increased sgACC coupling with non-DMN systems, most notably the Cognitive Control network, during and following sonication.`'